\structsection{Заключение}

В ходе производственной практики была исследована задача Smart 3D Packing, представляющая собой прикладную постановку 3D Bin Packing для укладки коробов на паллету. В отличие от абстрактной геометрической задачи, рассматриваемая постановка учитывает ограничения, характерные для складской логистики: габариты паллеты, грузоподъемность, допустимые повороты, опору, хрупкость, штабелируемость и время расчета.

В ходе работы были изучены условие задачи Smart 3D Packing, материалы НИР и курсовой работы, а также исходный код проекта. По результатам анализа восстановлен pipeline решения: генерация или загрузка JSON-датасета, запуск одного из решателей, валидация жестких ограничений, расчет метрик, сохранение результатов и визуализация укладки.

Были рассмотрены реализованные методы: greedy на основе extreme points и gravity drop, послойный MaxRects-решатель, Large Neighborhood Search, гибрид GAN+GA, branch-and-bound для малых задач и многопаллетное расширение. Для каждого метода описаны основная идея, роль в проекте, преимущества и ограничения.

Экспериментальная часть основана на сохраненных результатах репозитория. Сравнение на общих single-pallet наборах показало, что LNS является наиболее эффективным методом по итоговой метрике: 0.7111 на наборе из 100 задач, 0.7434 на наборе из 200 задач и 0.7298 на контрольном наборе. Его преимущество связано с тем, что метод умеет перестраивать уже найденную укладку, исправляя локально неудачные ранние решения greedy.

Greedy показал высокую скорость и идеальный временной балл, но уступил по утилизации объема и покрытию заказа. MaxRects достиг высокой плотности, но потерял качество по хрупкости. GAN+GA подтвердил перспективность гибридного ML-ориентированного поиска, однако в текущей реализации не превзошел LNS. Brute force оказался полезен для малых задач, но плохо масштабируется на более крупные экземпляры.

Таким образом, цель работы достигнута: выполнен анализ предметной области, формализована постановка, изучена реализация проекта и проведено сравнение методов. Наиболее рациональным методом для текущей single-pallet постановки является LNS; дальнейшее развитие связано с усилением destroy/repair-операторов и улучшением модели устойчивости.
